% Course: 77972 -- Physical Models of Biological Systems
% Academic year: 2026--2027 | Semester: A
% Actual lecture: 4 | Document chapter: 4
% Lecture date: 2026-11-10
% Title: קצבי מוטציה ספונטנית ומדדי מגוון ביולוגי
% Lecturer: Dr. Amir Erez
% Summarizer: Unknown (AI-assisted preparation)
% Prepared: 2026-11-12
% Editorial choice: generations g=1,...,G, synchronous binary fission, no death; the estimator uses K=2^G-1
%   consistently. Supplementary complements are marked as such in the text.

\chapter{קצבי מוטציה ספונטנית ומדדי מגוון ביולוגי}\label{chap:he04}
\noindent\textbf{הרצאה 4 \enspace|\enspace 10 בנובמבר 2026}

בהרצאה זו שתי שרשראות חשיבה אשר נראות רחוקות זו מזו, אך שתיהן עוסקות באותה שאלה: כיצד להוציא מספר משמעותי מהתפלגות. בחלק הראשון משלימים את ניתוח ניסוי התנודות של לוריא--דלברוק ומגזירים את אומד הנראות המרבית לקצב מוטציה ספונטנית. בחלק השני עוברים למגוון ביולוגי: אנטרופיית שנון, מספרי היל ואנטרופיית רניי, ולבסוף הכללה שמטילה משקל על דמיון בין מינים.

\subsection*{מטרות הלמידה}
\begin{itemize}
\item למדל גידול חיידקים סינכרוני ומוטציה ספונטנית כתהליך פואסוני של אירועים בודדים.
\item לגזור ולפרש את פונקציית האבות של מספר המוטנטים בתבנית גידול אחת.
\item לקבל אומד נראות מרבית לקצב המוטציה מתוך שיעור התבניות ללא מוטנט.
\item להסביר מדוע השונות של מספר המוטנטים גדולה מהתוחלת הפואסונית, ומדוע דרושות דווקא תבניות רבות.
\item להגדיר את אנטרופיית שנון, את מספרי היל ואת אנטרופיית רניי, ולפרש את "המספר האפקטיבי של מינים".
\item להכליל את המגוון באמצעות מטריצת דמיון בין מינים ומרחקים פילוגנטיים.
\end{itemize}

\section{חלק א: קצבי מוטציה ספונטנית}

\subsection*{המערך הניסיוני וההשערות המתחרות}
בניסוי מגדלים אוכלוסיית חיידקים ממקור קטן לאורך \(G\) דורות, בלי גורם ברירה. בדור \(G\) חושפים את התבניות לנגיף קטלני (או לאנטיביוטיקה), ורק תאים הנושאים מוטציה של עמידות שורדים. היסטורית התחרו שתי השערות:
\begin{itemize}
\item \emph{למארקיאנית / מושרית:} המוטציות מופיעות בתגובה לאתגר, ולכן כל תבנית אמורה להכיל בערך אותו מספר קטן של מוטנטים.
\item \emph{דרוויניסטית / ספונטנית:} המוטציות מתרחשות באקראי בזמן הגידול, לפני הברירה. מכיוון שמוטציה מוקדמת משאירה צאצאים רבים, מספר המוטנטים משתנה דרמטית בין תבניות לתבניות (תופעת ה"ג'קפוט").
\end{itemize}
המסקנה של לוריא--דלברוק היא שהנתונים אינם מתאימים למודל פואסוני פשוט של מוטציה מושרית: השונות גדולה מהתוחלת בסדרי גודל. התמליל גם מדגיש שהניסוי דורש רק זהירות רגילה (מקורות זהים, תנאי גידול זהים, אותו מספר דורות) ושקצבי מוטציה ספונטניים ניתנים לאומדן מספירה פשוטה מאוד.

\subsection*{מודל הגידול והסימונים}
יהי דור \(g=0,1,\dots,G\) מכיל \(2^{g}\) תאים (חלוקה בינארית סינכרונית, ללא מוות). מוטציה יכולה להתרחש בכל חלוקה. יהי \(\mu\) ההסתברות שמוטציה מתרחשת בחלוקה נתונה. מכיוון שהמוטציות נדירות, מספר אירועי המוטציה בדור \(g\) ממודל כפואסוני עם תוחלת
\begin{equation}\label{eq:he-bio-lambda-g}
\lambda_g = \mu\,2^{g-1}, \qquad g=1,\dots,G,
\end{equation}
כי בדור \(g\) מתרחשות \(2^{g-1}\) חלוקות. מוטציה המתרחשת בדור \(g\) חווה \(G-g\) דורות נוספים של התרבות, ולכן תורמת
\begin{equation}\label{eq:he-bio-clone-size}
a_g = 2^{G-g}
\end{equation}
תאי מוטנט לאוכלוסייה הסופית. מספר אירועי המוטציה הכולל בתבנית אחת הוא המשתנה המקרי \(M=\sum_{g=1}^{G}N_g\), ותוחלתו מסומנת באות גדולה:
\begin{equation}\label{eq:he-bio-Lambda-mean}
\Lambda\equiv\mathbb{E}[M]=\sum_{g=1}^{G}\lambda_g=\mu K,\qquad
K=\sum_{g=1}^{G}2^{g-1}=2^{G}-1.
\end{equation}
כאן \(K\) הוא מספר ההזדמנויות לחלוקה, ו-\(\Lambda\) (ולא \(M\)) הוא מספר האירועים הדטרמיניסטי הממוצע. מכיוון שהאוכלוסייה הסופית היא \(N=2^{G}\), מתקיים \(K=N-1\approx N\); האומד \(\eqref{eq:he-bio-mle}\) חייב להשתמש ב-\(K\) המתאים למוסכמה המוצהרת כאן.

\subsection*{פונקציית האבות של מספר המוטנטים}
יהי \(X\) מספר תאי המוטנט באוכלוסייה אחת לאחר הברירה. בהינתן אירועי המוטציה, \(X=\sum_{g}N_ga_g\), כאשר \(N_g\sim\text{Poisson}(\lambda_g)\) סופרים את אירועי המוטציה בדור \(g\). פונקציית ההסתברות היוצרת היא
\begin{align}
G_X(s)&=\mathbb{E}[s^{X}]
       =\prod_{g=1}^{G}\exp\!\bigl(\lambda_g(s^{a_g}-1)\bigr) \notag\\
      &=\exp\!\left(\sum_{g=1}^{G}\lambda_g(s^{a_g}-1)\right).
\label{eq:he-bio-gf}
\end{align}
מכאן שתי תוצאות מיידיות. ראשית, ההסתברות שאין מוטנט כלל בתבנית:
\begin{equation}\label{eq:he-bio-p0}
\Pr(X=0)=G_X(0)=\exp\!\left(-\sum_{g=1}^{G}\lambda_g\right)=e^{-\mu K}=e^{-\Lambda}.
\end{equation}
שנית, התוחלת של מספר המוטנטים:
\begin{equation}\label{eq:he-bio-mean}
\mathbb{E}[X]=G_X'(1)=\sum_{g=1}^{G}\lambda_g a_g
              =\mu\sum_{g=1}^{G}2^{g-1}2^{G-g}
              =\mu\,G\,2^{G-1}.
\end{equation}
מכיוון שה-\(N_g\) בלתי-תלויים, השונות של מספר המוטנטים מדויקת:
\begin{theorem}{שונות מספר המוטנטים וגורם פאנו}{variance-fano-he}
\begin{align}
\operatorname{Var}(X)
&=\sum_{g=1}^{G}a_g^{\,2}\operatorname{Var}(N_g)
 =\sum_{g=1}^{G}\lambda_g a_g^{\,2} \notag\\
&=\mu\sum_{g=1}^{G}2^{g-1}\bigl(2^{G-g}\bigr)^2
 =\mu\sum_{g=1}^{G}2^{2G-g-1}
 =\mu\,2^{G-1}(2^{G}-1),
\label{eq:he-bio-variance}
\end{align}
ולכן גורם פאנו הוא
\begin{equation}\label{eq:he-bio-var-ratio}
\boxed{\frac{\operatorname{Var}(X)}{\mathbb{E}[X]}=\frac{2^{G}-1}{G}} .
\end{equation}
\end{theorem}
\begin{proof}
מהעצמאות של \(N_g\): שונות של סכום בלתי-תלוי היא סכום השונויות, והשונות של כל \(a_gN_g\) היא \(a_g^2\operatorname{Var}(N_g) = a_g^2\lambda_g\). הסכום \(\sum_{g}2^{g-1}(2^{G-g})^2 = \sum_g 2^{2G-g-1} = 2^{2G-1}\sum_g 2^{-g} = 2^{2G-1}(1 - 2^{-G}) = 2^{G-1}(2^G-1)\). חלוקה בתוחלת \(\mu G 2^{G-1}\) נותנת את היחס.
\end{proof}
\begin{example}{בדיקה מספרית}{fano-check-he}
עבור \(G=10\) הדורות, גורם פאנו הוא \((1023)/10 \approx 102\). כך השונות גדולה מהתוחלת פי שני סדרי גודל לערך, וזו החתימה הכמותית של ה"ג'קפוטים": מוטציה מוקדמת אחת מוגברת על ידי הגידול האקספוננציאלי, ולכן ההתפלגות מוטה בחריפות ודגימות סופיות אינן מייצגות היטב את האירועים המוקדמים הנדירים.
\end{example}
\begin{remark}
\textbf{שונות מול מודל פואסוני נאיבי}\par
אילו המוטציות היו מושרות על ידי האתגר (למארקיאני), מספר המוטנטים היה בערך פואסוני, עם שונות השווה לתוחלת. השונות העל-פואסונית \(\eqref{eq:he-bio-var-ratio}\) שוללת זאת ומזהה את התהליך כמוטציה ספונטנית המוגברת בגידול אקספוננציאלי. הערה מהתמליל: האירועים המוקדמים היעילים כל כך נדירים, שעם מספר תבניות סופי המוטציה הנצפית הראשונה מתרחשת לרוב בדור מאוחר יחסית; הארכת הניסוי מחלישה את ההטיה, אך אינה מבטלת את השונות.
\end{remark}

\subsection*{אומד הנראות המרבית של \(\mu\)}
מבצעים \(C\) תבניות בלתי-תלויות, ו-\(C_0\) מהן ללא מוטנטים. הנראות היא בינומית:
\begin{equation}\label{eq:he-bio-likelihood}
\mathcal{L}(\mu)=\binom{C}{C_0}\bigl(e^{-\mu K}\bigr)^{C_0}
                  \bigl(1-e^{-\mu K}\bigr)^{C-C_0}.
\end{equation}
\begin{theorem}{אומד הנראות המרבית לקצב מוטציה}{mle-he}
\begin{equation}\label{eq:he-bio-mle}
\boxed{\widehat{\mu}=-\frac{1}{K}\ln\!\left(\frac{C_0}{C}\right)} .
\end{equation}
\end{theorem}
\begin{proof}
מגזירים את הלוג-נראות
\(\ell(\mu) = C_0(-\mu K) + (C-C_0)\ln(1-e^{-\mu K}) + \text{const}\) לפי \(\mu\) ומשווים לאפס:
\begin{equation}
\frac{\partial\ell}{\partial\mu} = -C_0K + (C-C_0)\frac{K e^{-\mu K}}{1-e^{-\mu K}} = 0
\;\Longrightarrow\; \frac{C-C_0}{1-e^{-\mu K}} = \frac{C_0}{e^{-\mu K}} .
\end{equation}
מכאן \(e^{-\mu K} = C_0/C\), והאומד \(\eqref{eq:he-bio-mle}\) נובע. הנגזרת השנייה שלילית, ולכן זהו מקסימום. \emph{מקרי גבול:} עבור \(C_0 = C\) (אין מוטטים כלל) האומד המתקבל \(0\), כלומר המקסימום נמצא על גבול הפרמטר; לא ניתן לאמוד קצב חיובי. עבור \(C_0 = 0\) יש \(\ln 0 = -\infty\) ואין אומד סופי.
\end{proof}
\begin{example}{אומדן קצב מוטציה בחיידקים}{mle-ex-he}
נניח \(C=100\) תבניות, ו-\(C_0=37\) מתוכן ללא מוטנטים עמידים. עם \(K=2^{G}-1\) הזדמנויות חלוקה מתקבל \(\widehat{\mu}=-\ln(0.37)/K \approx 1.0/K\). האומד תלוי רק בשיעור התבניות ה"סטריליות" ולא בספירת המוטנטים בג'קפוטים, וזה בדיוק מה שהופך את השיטה לחוסנת מפני השונות העצומה של \(X\).
\end{example}

\section{חלק ב: מגוון ביולוגי ואנטרופיה}

\subsection*{המערך ואנטרופיית שנון}
נתבונן בקהילה (או במערכת אקולוגית) עם \(m\) מינים (או סוגים) ששפעויותיהם היחסיות \(p_i>0\), \(\sum_{i=1}^{m}p_i=1\). אנטרופיית שנון היא
\begin{equation}\label{eq:he-bio-shannon}
H=-\sum_{i=1}^{m}p_i\ln p_i \quad (\text{nats}),
\end{equation}
והמספר האפקטיבי של מינים, הנקרא גם מספר היל מסדר \(1\), הוא
\begin{equation}\label{eq:he-bio-D1}
D_1=e^{H}.
\end{equation}
הפרשנות: \(D_1\) הוא מספר המינים השווים-בשפעות שהיו מניבים את אותה אנטרופיה. לדוגמה, עם שני מינים בשיעור \(p_1=p_2=1/2\):
\begin{equation}\label{eq:he-bio-two-species}
H=\ln 2,\qquad D_1=2.
\end{equation}
אם מין אחד שולט (\(p_1\approx 1\)), אזי \(H\approx 0\) ו-\(D_1\approx 1\). כך \(D_1\) מצמצם התפלגות שפעויות מוטה למספר מובן: "כמה מינים נפוצים באמת קיימים".

\subsection*{מספרי היל ואנטרופיית רניי}
מספרי היל מהווים משפחה חד-פרמטרית:
\begin{equation}\label{eq:he-bio-hill}
D_q=\left(\sum_{i=1}^{m}p_i^{\,q}\right)^{\!1/(1-q)},\qquad q\ge 0,\; q\neq 1,
\end{equation}
כאשר \(D_1\) מוגדר על ידי הרציפות כ-\(e^{H}\). בשקילות, אנטרופיית רניי מסדר \(q\) היא
\begin{equation}\label{eq:he-bio-renyi}
H_q=\frac{1}{1-q}\ln\!\left(\sum_{i=1}^{m}p_i^{\,q}\right),
\qquad D_q=e^{H_q}.
\end{equation}
הגבולות החשובים (כפי שניתנו בהרצאה):
\begin{itemize}
\item \(q\to 0\): \(D_0=m\), עושר המינים (כל המינים במשקל שווה, כולל נדירים).
\item \(q=1\): \(D_1=e^{H}\), המספר האפקטיבי של שנון.
\item \(q=2\): \(D_2=1/\sum p_i^2\), ההופכי של מדד סימפסון (ריכוז).
\item \(q\to\infty\): \(D_\infty=1/p_{\max}\), נשלט על ידי המין הנפוץ ביותר.
\end{itemize}
ככל ש-\(q\) גדל, המינים הנדירים מקבלים משקל קטן יותר; המשפחה נותנת ספקטרום ולא מספר יחיד. פרט להתפלגות שווה מנוונת, \(D_q\) יורדת ממשית ב-\(q\).
\begin{remark}
\textbf{עושר המינים שברירי}\par
\(D_0=m\) סופר כל מין, כולל כאלה נדירים מכדי להימדד בדגימה סופית או מתחת לגבול הזיהוי של המכשיר. התמליל מדגיש שעושר המינים לכן בלתי-יציב: העברת אתר הדגימה במטר, או החלפת הגלאי, יכולה לשנות את \(m\) במידה ניכרת. ערכי \(q\) גדולים נותנים מדדים יציבים יותר, כי המינים הנפוצים שולטים.
\end{remark}

\subsection*{מדד מגוון רגיש לדמיון (Leinster--Cobbold)}
עד כה התייחסנו לכל המינים כאל שונים לחלוטין. במציאות מינים שונים במידת הדמיון ביניהם. תהא \(Z\) מטריצת דמיון \(m\times m\) עם \(Z_{ii}=1\) ו-\(0\le Z_{ij}\le 1\) (למשל \(Z_{ij}=e^{-d_{ij}}\) עבור מרחק \(d_{ij}\), או דמיון טקסונומי). מגוון Leinster--Cobbold מסדר \(q\) הוא
\begin{equation}\label{eq:he-bio-lc}
{}^{q}D_Z(\mathbf p)=
\left(\sum_{i=1}^{m} p_i\left(\sum_{j=1}^{m}Z_{ij}p_j\right)^{\!q-1}\right)^{\!1/(1-q)},
\qquad q\neq 1,
\end{equation}
כאשר המקרה \(q=1\) נתון על ידי \(\exp\!\bigl(-\sum_i p_i\ln(\sum_j Z_{ij}p_j)\bigr)\). עבור \(Z=I\) מתקבלים מספרי היל \(\eqref{eq:he-bio-hill}\). הסכום הפנימי \(\sum_j Z_{ij}p_j\) הוא "השפעות האפקטיבית" של מין \(i\), כולל שכניו הדומים; הנוסחה נותנת משקל קטן יותר לאשכולות של מינים דומים.
\begin{theorem}{המקרה של שני מינים דומים באותה שפעות}{two-similar-he}
יהיו \(p_1=p_2=1/2\) ו-\(Z=\begin{psmallmatrix}1&z\\z&1\end{psmallmatrix}\), עם \(0\le z\le 1\). אזי לכל סדר \(q\),
\begin{equation}\label{eq:he-bio-two-similar}
\boxed{{}^{q}D_Z=\frac{2}{1+z}} .
\end{equation}
\end{theorem}
\begin{proof}
הסכום הפנימי זהה לשני המינים: \((Z\mathbf p)_1 = (Z\mathbf p)_2 = (1+z)/2\). לכן
\begin{equation}
\sum_i p_i (Z\mathbf p)_i^{q-1} = \frac12\left(\frac{1+z}{2}\right)^{q-1} + \frac12\left(\frac{1+z}{2}\right)^{q-1} = \left(\frac{1+z}{2}\right)^{q-1}.
\end{equation}
העלאה בחזקת \(1/(1-q)\) נותנת \((1+z)/2\) בהופכי, כלומר \(2/(1+z)\). עבור \(z=0\) (מינים שונים) מתקבל \(2\), ועבור \(z=1\) (מינים בלתי-ניתנים להבחנה) מתקבל \(1\). הגבול \(q\to 1\) של \((\frac{1+z}{2})^{q-1}\) מתכנס ל-1, ולכן הנוסחה ממשיכה כמו במקרה \(q=1\).
\end{proof}
\begin{remark}
\textbf{הניוון נובע משפעויות שוות}\par
שפעויות רגילות שוות הופכות את התוצאה לבלתי-תלויה ב-\(q\). זו התפרקות מקרית של הדוגמה, ולא כלל כללי: עבור שפעויות לא-שוות התוצאה תלויה ב-\(q\).
\end{remark}
\subsection*{מדד רייאו והגבול של סימפסון}
בוחרים שני פרטים באופן בלתי-תלוי מן הקהילה. אם אי-הדמיון בין המינים הוא \(d_{ij}=d_{ji}\), עם \(d_{ii}=0\), אי-הדמיון הממוצע ביניהם הוא מדד רייאו הריבועי:
\begin{equation}\label{eq:he-bio-rao}
Q=\sum_{i,j}p_i p_j d_{ij}.
\end{equation}
עבור מינים שונים לחלוטין, \(d_{ij}=1\) כאשר \(i\ne j\), ולכן
\begin{equation}\label{eq:he-bio-rao-simpson}
Q=\sum_{i\ne j}p_i p_j=1-\sum_i p_i^2=1-\frac{1}{D_2}.
\end{equation}
כלומר \(Q\) מצטמצם לגיוון גיני-סימפסון, ולא למספר האפקטיבי ההופכי של סימפסון \(D_2\) עצמו. אם \(Z_{ij}=1-d_{ij}\), אזי \(\sum_{ij}p_iZ_{ij}p_j=1-Q\), והמספר האפקטיבי הרגיש לדמיון מסדר שתיים הוא \(1/(1-Q)\). זה מבדיל בין אי-דמיון צפוי לבין ספירה אפקטיבית.

\subsection*{מרחקים פילוגנטיים}
עבור מערכות חיות הדמיון הרלוונטי נובע מעץ החיים. כל היצורים יוצאים מאב קדמון משותף יחיד; מרחקים אבולוציוניים זוגיים ניתנים לקריאה מעץ פילוגנטי (למשל מרצפים גנטיים). לשתי קהילות עשויים להיות אותו עושר מינים ואותה אנטרופיית שנון, אך אם אחת מורכבת ממינים קרובים גנטית והשנייה ממינים רחוקים, השנייה מגוונת יותר בכל מובן רגיש-לדמיון. התמליל ממחיש זאת במדגם מן הביצה, שבו שני מינים נראים שונים אך שייכים לאותה משפחה.
דמיון טקסונומי גס (ששימש בטבלת הדוגמה בהרצאה) הוא
\begin{equation}\label{eq:he-bio-taxsim}
Z_{ij}=
\begin{cases}
1, & \text{אותו מין},\\
0.8,& \text{אותו סוג, מינים שונים},\\
0.5,& \text{אותה משפחה, סוגים שונים},\\
0,  & \text{אחרת}.
\end{cases}
\end{equation}
מרחקים גנטיים מדויקים יותר אך דורשים ריצוף; התמליל מציין שניהול עצים פילוגנטיים הוא בעצמו תחום מחקר פעיל ושאין שיטה נכונה יחידה.
\begin{example}{פרפרים וחיידקי מעי}{diversity-ex-he}
ההרצאה משווה שתי מערכות אקולוגיות באמצעות הספקטרום הרציף \({}^{q}D_Z\):
\begin{itemize}
\item \emph{פרפרים בקנופי של יער גשם:} בבית גידול אחד המינים מרובים יותר אך רבים מהם מאותו סוג; בבית אחר יש פחות מינים, אך הם פרושים על פני סוגים רבים יותר. ספירת מינים בלבד מעדיפה את הראשון; שקלול לפי דמיון יכול להפוך את המסקנה עבור \(q\) גדולים.
\item \emph{מיקרוביום המעי:} לאם רזה ולאם שמנה קהילות חיידקים בעלות עושר נאיבי דומה. שקלול לפי דמיון רצף מוריד את המגוון האפקטיבי הרבה יותר אצל האם השמנה, כי הקהילה שלה נשלטת על ידי זנים קרובים רבים; המדד הנאיבי (ללא שקלול) מגזים את המגוון שלה.
\end{itemize}
המסקנה של התמליל: מספרי היל ללא שקלול עלולים להטעות בהחלטות שימור או רפואיות, והספקטרום המשוקלל בדמיון אינפורמטיבי יותר.
\end{example}

\section*{שאלות חזרה ופתרונות}
\begin{enumerate}
\item \emph{מדוע השונות של \(X\) גדולה בהרבה מתוחלתו?} כי כל מוטציה תורמת \(2^{G-g}\) תאים, והערך גדל אקספוננציאלית ככל שהמוטציה מוקדמת יותר. גורם פאנו הוא \((2^G-1)/G\).
\item \emph{מהו \(\Pr(X=0)\) ומדוע הוא שימושי?} \(\Pr(X=0)=e^{-\mu K}\). הוא נותן אומדן לקצב המוטציה מהשיעור \(C_0/C\) בלי לספור מוטנטים כלל, ולכן חוסן מאבדן ג'קפוטים.
\item \emph{מדוע \(D_0\) פחות יציב מ-\(D_2\)?} עושר המינים סופר גם מינים נדירים מאוד, אשר הופעתם תלויה בגודל המדגם ובגבול הזיהוי. \(D_2\) נשלט על ידי המינים הנפוצים.
\item \emph{מהו ההבדל בין \(Q\) לבין \(1-1/D_2\)?} אין הבדל כשכל המינים שונים; \(Q\) הוא אי-דמיון צפוי, ואילו \(1-1/D_2\) הוא ההשלכה עבור מינים שונים לחלוטין.
\item \emph{מדוע עבור שני מינים באותה שפעות התוצאה בלתי-תלויה ב-\(q\)?} כי אז \((Z\mathbf p)_1=(Z\mathbf p)_2\) מפורק את הביטוי, והתלות ב-\(q\) נעלמת.
\end{enumerate}

\subsection*{קריאה נוספת}
\begin{itemize}
\item \textenglish{Philip Nelson}, \emph{\textenglish{Physical Models of Living Systems}}, פרקים על תהליכים סטוכסטיים, מוטציות מידע ואנטרופיה.
\item \textenglish{Luria, S. E. and Delbrück, M. (1943), ``Mutations of bacteria from virus sensitivity to virus resistance'', \emph{Genetics} 28, 491--511.}
\item \textenglish{Hill, M. O. (1973), ``Diversity and evenness: a unifying notation and its consequences'', \emph{Ecology} 54, 427--432.}
\item \textenglish{Rényi, A. (1961), ``On measures of entropy and information'', \emph{Proc. 4th Berkeley Symp.} 1, 547--561.}
\item \textenglish{Leinster, T. and Cobbold, C. A. (2012), ``Measuring diversity: the importance of species similarity'', \emph{Ecology} 93, 477--489.}
\end{itemize}
